% --- Archon physics formalization source begin ---
% archon:physics
% archon:covers PhyXMiniProblems/problem_phyx_mini_0884.lean
% archon:source-report reports/phyx_mini/problem_phyx_mini_0884.source.json

\chapter{Physics problem phyx\_mini\_0884}
\label{ch:PhyXMiniProblems_problem_phyx_mini_0884}

\paragraph{Problem source.}
\#\# Physical scenario

A closed, square loop is formed with 40 cm of wire having \textbackslash{}( R = 0.10 \textbackslash{}, \textbackslash{}Omega \textbackslash{}), as shown in figure. A 0.50 T magnetic field is perpendicular to the loop. At \textbackslash{}( t = 0 \textbackslash{}, \textbackslash{}text\{s\} \textbackslash{}), two diagonally opposite corners of the loop begin to move apart at 0.293 m/s.

\#\# Question

How long does it take the loop to collapse to a straight line?

\#\# Answer choices

- A: A: \textbackslash{}( 0.7 \textbackslash{}

- B: B: \textbackslash{}( 0.3 \textbackslash{}

- C: C: \textbackslash{}( 0.1 \textbackslash{}

- D: D: \textbackslash{}( 0.9 \textbackslash{}

\#\# Figure caption (auxiliary; use the image as primary evidence)

The image depicts a diagram with a rhombus at its center surrounded by a pattern of blue dots scattered uniformly across the background. A border is drawn around the rhombus, highlighted in light brown. Two green arrows are positioned horizontally at the midpoint of the rhombus, pointing in opposite directions—left and right. Each arrow is labeled with the text "0.293 m/s" indicating their velocities.

\#\# Dataset metadata

Electromagnetic Induction; Physical Model Grounding Reasoning, Spatial Relation Reasoning

\paragraph{Recorded answer/context.}
C: C: \textbackslash{}( 0.1 \textbackslash{}

\paragraph{Figure/image path.}
/root/proposal\_for\_physic/science-mango/phyx\_mini\_run/phyx\_data/test\_image/884.png

\paragraph{Formalization target.}
create a compiling Lean file with sorry bodies at `PhyXMiniProblems/problem\_phyx\_mini\_0884.lean`. The Lean declarations must preserve the physical quantities, dimensions or dimensional roles, figure labels, governing-law hypotheses, and final relation expressed by this problem.
Use Mathlib/Physlib names found through LeanExplore where available. If a physics API is missing, introduce faithful local abstractions rather than scalar placeholder aliases.

\begin{theorem}[Physics formalization target]
\label{thm:physics:phyx_mini_0884:target}
\lean{PhyXMiniProblems.ProblemPhyXMini0884.problem_phyx_mini_0884}
\uses{def:physics:phyx-mini-0884:phyxminiproblems-problemphyxmini0884-collapsingsquareloopsetup, def:physics:phyx-mini-0884:phyxminiproblems-problemphyxmini0884-collapseelapsedtimeinseconds, def:physics:phyx-mini-0884:phyxminiproblems-problemphyxmini0884-matchescollapsingsquareloopdescription, def:physics:phyx-mini-0884:phyxminiproblems-problemphyxmini0884-matchesprimarycollapsingloopfigure, def:physics:phyx-mini-0884:phyxminiproblems-problemphyxmini0884-hasphysicalcollapsingloopparameters, def:physics:phyx-mini-0884:phyxminiproblems-problemphyxmini0884-satisfiesinextensiblesquaregeometry, def:physics:phyx-mini-0884:phyxminiproblems-problemphyxmini0884-satisfiesoppositecornerseparationkinematics, def:physics:phyx-mini-0884:phyxminiproblems-problemphyxmini0884-isfirststraightlinecollapse, def:physics:phyx-mini-0884:phyxminiproblems-problemphyxmini0884-recordeddatasetanswer, def:physics:phyx-mini-0884:phyxminiproblems-problemphyxmini0884-roundstonearesttenthsecond, def:physics:phyx-mini-0884:phyxminiproblems-problemphyxmini0884-answermatchescollapsetime}
The assigned autoformalize agent should translate this physics problem into Lean declarations in the covered file, with theorem and lemma proof bodies written as `by sorry`.
\end{theorem}
\begin{proof}
This is an autoformalization task, not a proof task. Produce statements that can later be proved by the physics prover without weakening the physical model.
\end{proof}

% --- Archon declaration topology begin ---
\section{Declaration topology}
\begin{definition}[speed Dimension]
  \label{def:physics:phyx-mini-0884:phyxminiproblems-problemphyxmini0884-speeddimension}
  \lean{PhyXMiniProblems.ProblemPhyXMini0884.speedDimension}
  The physical dimension L T⁻¹ of speed
\end{definition}
\begin{proof}
  This is the declared model definition of speed Dimension.
\end{proof}
\begin{definition}[electrical Resistance Dimension]
  \label{def:physics:phyx-mini-0884:phyxminiproblems-problemphyxmini0884-electricalresistancedimension}
  \lean{PhyXMiniProblems.ProblemPhyXMini0884.electricalResistanceDimension}
  The physical dimension M L² T⁻¹ C⁻² of electrical resistance
\end{definition}
\begin{proof}
  This is the declared model definition of electrical Resistance Dimension.
\end{proof}
\begin{definition}[magnetic Flux Density Dimension]
  \label{def:physics:phyx-mini-0884:phyxminiproblems-problemphyxmini0884-magneticfluxdensitydimension}
  \lean{PhyXMiniProblems.ProblemPhyXMini0884.magneticFluxDensityDimension}
  The physical dimension M T⁻¹ C⁻¹ of magnetic flux density
\end{definition}
\begin{proof}
  This is the declared model definition of magnetic Flux Density Dimension.
\end{proof}
\begin{definition}[Length Magnitude]
  \label{def:physics:phyx-mini-0884:phyxminiproblems-problemphyxmini0884-lengthmagnitude}
  \lean{PhyXMiniProblems.ProblemPhyXMini0884.LengthMagnitude}
  A nonnegative, unit-independent physical length
\end{definition}
\begin{proof}
  This is the declared model definition of Length Magnitude.
\end{proof}
\begin{definition}[Time Magnitude]
  \label{def:physics:phyx-mini-0884:phyxminiproblems-problemphyxmini0884-timemagnitude}
  \lean{PhyXMiniProblems.ProblemPhyXMini0884.TimeMagnitude}
  A nonnegative, unit-independent physical time
\end{definition}
\begin{proof}
  This is the declared model definition of Time Magnitude.
\end{proof}
\begin{definition}[Speed Magnitude]
  \label{def:physics:phyx-mini-0884:phyxminiproblems-problemphyxmini0884-speedmagnitude}
  \lean{PhyXMiniProblems.ProblemPhyXMini0884.SpeedMagnitude}
  \uses{def:physics:phyx-mini-0884:phyxminiproblems-problemphyxmini0884-speeddimension}
  A nonnegative, unit-independent physical speed
\end{definition}
\begin{proof}
  This is the declared model definition of Speed Magnitude.
\end{proof}
\begin{definition}[Resistance Magnitude]
  \label{def:physics:phyx-mini-0884:phyxminiproblems-problemphyxmini0884-resistancemagnitude}
  \lean{PhyXMiniProblems.ProblemPhyXMini0884.ResistanceMagnitude}
  \uses{def:physics:phyx-mini-0884:phyxminiproblems-problemphyxmini0884-electricalresistancedimension}
  A nonnegative, unit-independent electrical resistance
\end{definition}
\begin{proof}
  This is the declared model definition of Resistance Magnitude.
\end{proof}
\begin{definition}[Magnetic Flux Density Magnitude]
  \label{def:physics:phyx-mini-0884:phyxminiproblems-problemphyxmini0884-magneticfluxdensitymagnitude}
  \lean{PhyXMiniProblems.ProblemPhyXMini0884.MagneticFluxDensityMagnitude}
  \uses{def:physics:phyx-mini-0884:phyxminiproblems-problemphyxmini0884-magneticfluxdensitydimension}
  A nonnegative, unit-independent magnetic-flux-density magnitude
\end{definition}
\begin{proof}
  This is the declared model definition of Magnetic Flux Density Magnitude.
\end{proof}
\begin{definition}[length Readout]
  \label{def:physics:phyx-mini-0884:phyxminiproblems-problemphyxmini0884-lengthreadout}
  \lean{PhyXMiniProblems.ProblemPhyXMini0884.lengthReadout}
  \uses{def:physics:phyx-mini-0884:phyxminiproblems-problemphyxmini0884-lengthmagnitude}
  Read a physical length in the coherent length unit selected by units
\end{definition}
\begin{proof}
  This is the declared model definition of length Readout.
\end{proof}
\begin{definition}[time Readout]
  \label{def:physics:phyx-mini-0884:phyxminiproblems-problemphyxmini0884-timereadout}
  \lean{PhyXMiniProblems.ProblemPhyXMini0884.timeReadout}
  \uses{def:physics:phyx-mini-0884:phyxminiproblems-problemphyxmini0884-timemagnitude}
  Read a physical time in the coherent time unit selected by units
\end{definition}
\begin{proof}
  This is the declared model definition of time Readout.
\end{proof}
\begin{definition}[speed Readout]
  \label{def:physics:phyx-mini-0884:phyxminiproblems-problemphyxmini0884-speedreadout}
  \lean{PhyXMiniProblems.ProblemPhyXMini0884.speedReadout}
  \uses{def:physics:phyx-mini-0884:phyxminiproblems-problemphyxmini0884-speedmagnitude}
  Read a physical speed in the coherent speed unit selected by units
\end{definition}
\begin{proof}
  This is the declared model definition of speed Readout.
\end{proof}
\begin{definition}[resistance Readout]
  \label{def:physics:phyx-mini-0884:phyxminiproblems-problemphyxmini0884-resistancereadout}
  \lean{PhyXMiniProblems.ProblemPhyXMini0884.resistanceReadout}
  \uses{def:physics:phyx-mini-0884:phyxminiproblems-problemphyxmini0884-resistancemagnitude}
  Read a resistance in the coherent resistance unit selected by units
\end{definition}
\begin{proof}
  This is the declared model definition of resistance Readout.
\end{proof}
\begin{definition}[magnetic Flux Density Readout]
  \label{def:physics:phyx-mini-0884:phyxminiproblems-problemphyxmini0884-magneticfluxdensityreadout}
  \lean{PhyXMiniProblems.ProblemPhyXMini0884.magneticFluxDensityReadout}
  \uses{def:physics:phyx-mini-0884:phyxminiproblems-problemphyxmini0884-magneticfluxdensitymagnitude}
  Read magnetic flux density in the coherent unit selected by units
\end{definition}
\begin{proof}
  This is the declared model definition of magnetic Flux Density Readout.
\end{proof}
\begin{definition}[length In Meters]
  \label{def:physics:phyx-mini-0884:phyxminiproblems-problemphyxmini0884-lengthinmeters}
  \lean{PhyXMiniProblems.ProblemPhyXMini0884.lengthInMeters}
  \uses{def:physics:phyx-mini-0884:phyxminiproblems-problemphyxmini0884-lengthmagnitude, def:physics:phyx-mini-0884:phyxminiproblems-problemphyxmini0884-lengthreadout}
  Coherent-SI metre readout of a physical length
\end{definition}
\begin{proof}
  This is the declared model definition of length In Meters.
\end{proof}
\begin{definition}[length In Centimeters]
  \label{def:physics:phyx-mini-0884:phyxminiproblems-problemphyxmini0884-lengthincentimeters}
  \lean{PhyXMiniProblems.ProblemPhyXMini0884.lengthInCentimeters}
  \uses{def:physics:phyx-mini-0884:phyxminiproblems-problemphyxmini0884-lengthmagnitude, def:physics:phyx-mini-0884:phyxminiproblems-problemphyxmini0884-lengthinmeters}
  Centimetre readout used for the stated total wire length
\end{definition}
\begin{proof}
  This is the declared model definition of length In Centimeters.
\end{proof}
\begin{definition}[time In Seconds]
  \label{def:physics:phyx-mini-0884:phyxminiproblems-problemphyxmini0884-timeinseconds}
  \lean{PhyXMiniProblems.ProblemPhyXMini0884.timeInSeconds}
  \uses{def:physics:phyx-mini-0884:phyxminiproblems-problemphyxmini0884-timemagnitude, def:physics:phyx-mini-0884:phyxminiproblems-problemphyxmini0884-timereadout}
  Coherent-SI second readout of a physical time
\end{definition}
\begin{proof}
  This is the declared model definition of time In Seconds.
\end{proof}
\begin{definition}[speed In Meters Per Second]
  \label{def:physics:phyx-mini-0884:phyxminiproblems-problemphyxmini0884-speedinmeterspersecond}
  \lean{PhyXMiniProblems.ProblemPhyXMini0884.speedInMetersPerSecond}
  \uses{def:physics:phyx-mini-0884:phyxminiproblems-problemphyxmini0884-speedmagnitude, def:physics:phyx-mini-0884:phyxminiproblems-problemphyxmini0884-speedreadout}
  Coherent-SI metre-per-second readout of a physical speed
\end{definition}
\begin{proof}
  This is the declared model definition of speed In Meters Per Second.
\end{proof}
\begin{definition}[resistance In Ohms]
  \label{def:physics:phyx-mini-0884:phyxminiproblems-problemphyxmini0884-resistanceinohms}
  \lean{PhyXMiniProblems.ProblemPhyXMini0884.resistanceInOhms}
  \uses{def:physics:phyx-mini-0884:phyxminiproblems-problemphyxmini0884-resistancemagnitude, def:physics:phyx-mini-0884:phyxminiproblems-problemphyxmini0884-resistancereadout}
  Coherent-SI ohm readout of a physical resistance
\end{definition}
\begin{proof}
  This is the declared model definition of resistance In Ohms.
\end{proof}
\begin{definition}[magnetic Flux Density In Teslas]
  \label{def:physics:phyx-mini-0884:phyxminiproblems-problemphyxmini0884-magneticfluxdensityinteslas}
  \lean{PhyXMiniProblems.ProblemPhyXMini0884.magneticFluxDensityInTeslas}
  \uses{def:physics:phyx-mini-0884:phyxminiproblems-problemphyxmini0884-magneticfluxdensitymagnitude, def:physics:phyx-mini-0884:phyxminiproblems-problemphyxmini0884-magneticfluxdensityreadout}
  Coherent-SI tesla readout of a magnetic-flux-density magnitude
\end{definition}
\begin{proof}
  This is the declared model definition of magnetic Flux Density In Teslas.
\end{proof}
\begin{definition}[Loop Corner]
  \label{def:physics:phyx-mini-0884:phyxminiproblems-problemphyxmini0884-loopcorner}
  \lean{PhyXMiniProblems.ProblemPhyXMini0884.LoopCorner}
  The four vertices of the diamond-oriented square in image 884.png
\end{definition}
\begin{proof}
  This is the declared model definition of Loop Corner.
\end{proof}
\begin{definition}[Wire Segment]
  \label{def:physics:phyx-mini-0884:phyxminiproblems-problemphyxmini0884-wiresegment}
  \lean{PhyXMiniProblems.ProblemPhyXMini0884.WireSegment}
  The four pieces of wire joining successive visible corners
\end{definition}
\begin{proof}
  This is the declared model definition of Wire Segment.
\end{proof}
\begin{definition}[Pulled Corner]
  \label{def:physics:phyx-mini-0884:phyxminiproblems-problemphyxmini0884-pulledcorner}
  \lean{PhyXMiniProblems.ProblemPhyXMini0884.PulledCorner}
  The two diagonally opposite corners to which the green arrows are attached
\end{definition}
\begin{proof}
  This is the declared model definition of Pulled Corner.
\end{proof}
\begin{definition}[Horizontal Direction]
  \label{def:physics:phyx-mini-0884:phyxminiproblems-problemphyxmini0884-horizontaldirection}
  \lean{PhyXMiniProblems.ProblemPhyXMini0884.HorizontalDirection}
  Horizontal directions of the two velocity arrows in the figure
\end{definition}
\begin{proof}
  This is the declared model definition of Horizontal Direction.
\end{proof}
\begin{definition}[Field Page Direction]
  \label{def:physics:phyx-mini-0884:phyxminiproblems-problemphyxmini0884-fieldpagedirection}
  \lean{PhyXMiniProblems.ProblemPhyXMini0884.FieldPageDirection}
  Direction represented by the conventional blue dot field markers
\end{definition}
\begin{proof}
  This is the declared model definition of Field Page Direction.
\end{proof}
\begin{definition}[Loop Topology]
  \label{def:physics:phyx-mini-0884:phyxminiproblems-problemphyxmini0884-looptopology}
  \lean{PhyXMiniProblems.ProblemPhyXMini0884.LoopTopology}
  Topology of the physical wire
\end{definition}
\begin{proof}
  This is the declared model definition of Loop Topology.
\end{proof}
\begin{definition}[Initial Loop Shape]
  \label{def:physics:phyx-mini-0884:phyxminiproblems-problemphyxmini0884-initialloopshape}
  \lean{PhyXMiniProblems.ProblemPhyXMini0884.InitialLoopShape}
  Initial shape asserted by the problem prose
\end{definition}
\begin{proof}
  This is the declared model definition of Initial Loop Shape.
\end{proof}
\begin{definition}[Wire Deformation Model]
  \label{def:physics:phyx-mini-0884:phyxminiproblems-problemphyxmini0884-wiredeformationmodel}
  \lean{PhyXMiniProblems.ProblemPhyXMini0884.WireDeformationModel}
  Mechanical idealization used while the square deforms through rhombi
\end{definition}
\begin{proof}
  This is the declared model definition of Wire Deformation Model.
\end{proof}
\begin{definition}[Field Loop Orientation]
  \label{def:physics:phyx-mini-0884:phyxminiproblems-problemphyxmini0884-fieldlooporientation}
  \lean{PhyXMiniProblems.ProblemPhyXMini0884.FieldLoopOrientation}
  Spatial relation between the magnetic field and the loop plane
\end{definition}
\begin{proof}
  This is the declared model definition of Field Loop Orientation.
\end{proof}
\begin{definition}[Field Uniformity]
  \label{def:physics:phyx-mini-0884:phyxminiproblems-problemphyxmini0884-fielduniformity}
  \lean{PhyXMiniProblems.ProblemPhyXMini0884.FieldUniformity}
  Uniformity idealization for the applied magnetic field
\end{definition}
\begin{proof}
  This is the declared model definition of Field Uniformity.
\end{proof}
\begin{definition}[expected Arrow Direction]
  \label{def:physics:phyx-mini-0884:phyxminiproblems-problemphyxmini0884-expectedarrowdirection}
  \lean{PhyXMiniProblems.ProblemPhyXMini0884.expectedArrowDirection}
  \uses{def:physics:phyx-mini-0884:phyxminiproblems-problemphyxmini0884-pulledcorner, def:physics:phyx-mini-0884:phyxminiproblems-problemphyxmini0884-horizontaldirection}
  The direction expected for each of the two outward velocity arrows
\end{definition}
\begin{proof}
  This is the declared model definition of expected Arrow Direction.
\end{proof}
\begin{definition}[Collapsing Loop Figure]
  \label{def:physics:phyx-mini-0884:phyxminiproblems-problemphyxmini0884-collapsingloopfigure}
  \lean{PhyXMiniProblems.ProblemPhyXMini0884.CollapsingLoopFigure}
  \uses{def:physics:phyx-mini-0884:phyxminiproblems-problemphyxmini0884-loopcorner, def:physics:phyx-mini-0884:phyxminiproblems-problemphyxmini0884-wiresegment, def:physics:phyx-mini-0884:phyxminiproblems-problemphyxmini0884-pulledcorner, def:physics:phyx-mini-0884:phyxminiproblems-problemphyxmini0884-horizontaldirection, def:physics:phyx-mini-0884:phyxminiproblems-problemphyxmini0884-fieldpagedirection}
  Literal information transcribed from the primary raster. Displayed speeds are scalar metre-per-second readouts; they are calibrated to independent physical speed quantities below
\end{definition}
\begin{proof}
  This is the declared model definition of Collapsing Loop Figure.
\end{proof}
\begin{definition}[Collapsing Square Loop Setup]
  \label{def:physics:phyx-mini-0884:phyxminiproblems-problemphyxmini0884-collapsingsquareloopsetup}
  \lean{PhyXMiniProblems.ProblemPhyXMini0884.CollapsingSquareLoopSetup}
  \uses{def:physics:phyx-mini-0884:phyxminiproblems-problemphyxmini0884-lengthmagnitude, def:physics:phyx-mini-0884:phyxminiproblems-problemphyxmini0884-timemagnitude, def:physics:phyx-mini-0884:phyxminiproblems-problemphyxmini0884-speedmagnitude, def:physics:phyx-mini-0884:phyxminiproblems-problemphyxmini0884-resistancemagnitude, def:physics:phyx-mini-0884:phyxminiproblems-problemphyxmini0884-magneticfluxdensitymagnitude, def:physics:phyx-mini-0884:phyxminiproblems-problemphyxmini0884-pulledcorner, def:physics:phyx-mini-0884:phyxminiproblems-problemphyxmini0884-looptopology, def:physics:phyx-mini-0884:phyxminiproblems-problemphyxmini0884-initialloopshape, def:physics:phyx-mini-0884:phyxminiproblems-problemphyxmini0884-wiredeformationmodel, def:physics:phyx-mini-0884:phyxminiproblems-problemphyxmini0884-fieldlooporientation, def:physics:phyx-mini-0884:phyxminiproblems-problemphyxmini0884-fielduniformity, def:physics:phyx-mini-0884:phyxminiproblems-problemphyxmini0884-collapsingloopfigure}
  Independent physical data for the deforming loop. In particular, collapseTime and pulledDiagonalLength are observables, not definitions made from 0.1 s or from any answer choice
\end{definition}
\begin{proof}
  This is the declared model definition of Collapsing Square Loop Setup.
\end{proof}
\begin{definition}[pulled Diagonal In Meters]
  \label{def:physics:phyx-mini-0884:phyxminiproblems-problemphyxmini0884-pulleddiagonalinmeters}
  \lean{PhyXMiniProblems.ProblemPhyXMini0884.pulledDiagonalInMeters}
  \uses{def:physics:phyx-mini-0884:phyxminiproblems-problemphyxmini0884-timemagnitude, def:physics:phyx-mini-0884:phyxminiproblems-problemphyxmini0884-lengthinmeters, def:physics:phyx-mini-0884:phyxminiproblems-problemphyxmini0884-collapsingsquareloopsetup}
  Length of the pulled left-right diagonal at a physical time, in metres
\end{definition}
\begin{proof}
  This is the declared model definition of pulled Diagonal In Meters.
\end{proof}
\begin{definition}[side Length In Meters]
  \label{def:physics:phyx-mini-0884:phyxminiproblems-problemphyxmini0884-sidelengthinmeters}
  \lean{PhyXMiniProblems.ProblemPhyXMini0884.sideLengthInMeters}
  \uses{def:physics:phyx-mini-0884:phyxminiproblems-problemphyxmini0884-lengthinmeters, def:physics:phyx-mini-0884:phyxminiproblems-problemphyxmini0884-collapsingsquareloopsetup}
  Side length of the inextensible rhombus, in metres
\end{definition}
\begin{proof}
  This is the declared model definition of side Length In Meters.
\end{proof}
\begin{definition}[collapse Elapsed Time In Seconds]
  \label{def:physics:phyx-mini-0884:phyxminiproblems-problemphyxmini0884-collapseelapsedtimeinseconds}
  \lean{PhyXMiniProblems.ProblemPhyXMini0884.collapseElapsedTimeInSeconds}
  \uses{def:physics:phyx-mini-0884:phyxminiproblems-problemphyxmini0884-timeinseconds, def:physics:phyx-mini-0884:phyxminiproblems-problemphyxmini0884-collapsingsquareloopsetup}
  Elapsed time from the onset of pulling to collapse, in seconds
\end{definition}
\begin{proof}
  This is the declared model definition of collapse Elapsed Time In Seconds.
\end{proof}
\begin{definition}[Matches Collapsing Square Loop Description]
  \label{def:physics:phyx-mini-0884:phyxminiproblems-problemphyxmini0884-matchescollapsingsquareloopdescription}
  \lean{PhyXMiniProblems.ProblemPhyXMini0884.MatchesCollapsingSquareLoopDescription}
  \uses{def:physics:phyx-mini-0884:phyxminiproblems-problemphyxmini0884-lengthincentimeters, def:physics:phyx-mini-0884:phyxminiproblems-problemphyxmini0884-timeinseconds, def:physics:phyx-mini-0884:phyxminiproblems-problemphyxmini0884-resistanceinohms, def:physics:phyx-mini-0884:phyxminiproblems-problemphyxmini0884-magneticfluxdensityinteslas, def:physics:phyx-mini-0884:phyxminiproblems-problemphyxmini0884-collapsingsquareloopsetup}
  Qualitative prose assumptions and the three non-figure numerical data
\end{definition}
\begin{proof}
  This is the declared model definition of Matches Collapsing Square Loop Description.
\end{proof}
\begin{definition}[Matches Primary Collapsing Loop Figure]
  \label{def:physics:phyx-mini-0884:phyxminiproblems-problemphyxmini0884-matchesprimarycollapsingloopfigure}
  \lean{PhyXMiniProblems.ProblemPhyXMini0884.MatchesPrimaryCollapsingLoopFigure}
  \uses{def:physics:phyx-mini-0884:phyxminiproblems-problemphyxmini0884-speedinmeterspersecond, def:physics:phyx-mini-0884:phyxminiproblems-problemphyxmini0884-expectedarrowdirection, def:physics:phyx-mini-0884:phyxminiproblems-problemphyxmini0884-collapsingsquareloopsetup}
  The primary-image facts, including the two separately drawn speeds. This structure contains no collapse-time value and no answer-choice data
\end{definition}
\begin{proof}
  This is the declared model definition of Matches Primary Collapsing Loop Figure.
\end{proof}
\begin{definition}[Has Physical Collapsing Loop Parameters]
  \label{def:physics:phyx-mini-0884:phyxminiproblems-problemphyxmini0884-hasphysicalcollapsingloopparameters}
  \lean{PhyXMiniProblems.ProblemPhyXMini0884.HasPhysicalCollapsingLoopParameters}
  \uses{def:physics:phyx-mini-0884:phyxminiproblems-problemphyxmini0884-lengthinmeters, def:physics:phyx-mini-0884:phyxminiproblems-problemphyxmini0884-timeinseconds, def:physics:phyx-mini-0884:phyxminiproblems-problemphyxmini0884-speedinmeterspersecond, def:physics:phyx-mini-0884:phyxminiproblems-problemphyxmini0884-resistanceinohms, def:physics:phyx-mini-0884:phyxminiproblems-problemphyxmini0884-magneticfluxdensityinteslas, def:physics:phyx-mini-0884:phyxminiproblems-problemphyxmini0884-collapsingsquareloopsetup, def:physics:phyx-mini-0884:phyxminiproblems-problemphyxmini0884-sidelengthinmeters}
  Positivity and chronological conditions selecting the physical branch
\end{definition}
\begin{proof}
  This is the declared model definition of Has Physical Collapsing Loop Parameters.
\end{proof}
\begin{definition}[Satisfies Inextensible Square Geometry]
  \label{def:physics:phyx-mini-0884:phyxminiproblems-problemphyxmini0884-satisfiesinextensiblesquaregeometry}
  \lean{PhyXMiniProblems.ProblemPhyXMini0884.SatisfiesInextensibleSquareGeometry}
  \uses{def:physics:phyx-mini-0884:phyxminiproblems-problemphyxmini0884-lengthreadout, def:physics:phyx-mini-0884:phyxminiproblems-problemphyxmini0884-collapsingsquareloopsetup}
  Fixed-perimeter and initial-square geometry, stated in every coherent unit system. The first equation says the 40 cm wire remains four equal sides; the second is the Pythagorean relation for the initial square diagonal
\end{definition}
\begin{proof}
  This is the declared model definition of Satisfies Inextensible Square Geometry.
\end{proof}
\begin{definition}[Satisfies Opposite Corner Separation Kinematics]
  \label{def:physics:phyx-mini-0884:phyxminiproblems-problemphyxmini0884-satisfiesoppositecornerseparationkinematics}
  \lean{PhyXMiniProblems.ProblemPhyXMini0884.SatisfiesOppositeCornerSeparationKinematics}
  \uses{def:physics:phyx-mini-0884:phyxminiproblems-problemphyxmini0884-timemagnitude, def:physics:phyx-mini-0884:phyxminiproblems-problemphyxmini0884-lengthreadout, def:physics:phyx-mini-0884:phyxminiproblems-problemphyxmini0884-timereadout, def:physics:phyx-mini-0884:phyxminiproblems-problemphyxmini0884-speedreadout, def:physics:phyx-mini-0884:phyxminiproblems-problemphyxmini0884-collapsingsquareloopsetup}
  The separation of the pulled corners grows at the sum of their two outward speeds. This is a general kinematic law on the interval from the start event through the independently modeled collapse event; it contains no requested time value
\end{definition}
\begin{proof}
  This is the declared model definition of Satisfies Opposite Corner Separation Kinematics.
\end{proof}
\begin{definition}[Is Straight Line Configuration]
  \label{def:physics:phyx-mini-0884:phyxminiproblems-problemphyxmini0884-isstraightlineconfiguration}
  \lean{PhyXMiniProblems.ProblemPhyXMini0884.IsStraightLineConfiguration}
  \uses{def:physics:phyx-mini-0884:phyxminiproblems-problemphyxmini0884-timemagnitude, def:physics:phyx-mini-0884:phyxminiproblems-problemphyxmini0884-lengthreadout, def:physics:phyx-mini-0884:phyxminiproblems-problemphyxmini0884-collapsingsquareloopsetup}
  For an inextensible four-sided loop pulled along one diagonal, the degenerate straight-line configuration occurs when that diagonal equals two side lengths. This geometric event definition has no numerical time in it
\end{definition}
\begin{proof}
  This is the declared model definition of Is Straight Line Configuration.
\end{proof}
\begin{definition}[Is First Straight Line Collapse]
  \label{def:physics:phyx-mini-0884:phyxminiproblems-problemphyxmini0884-isfirststraightlinecollapse}
  \lean{PhyXMiniProblems.ProblemPhyXMini0884.IsFirstStraightLineCollapse}
  \uses{def:physics:phyx-mini-0884:phyxminiproblems-problemphyxmini0884-timeinseconds, def:physics:phyx-mini-0884:phyxminiproblems-problemphyxmini0884-collapsingsquareloopsetup, def:physics:phyx-mini-0884:phyxminiproblems-problemphyxmini0884-isstraightlineconfiguration}
  The modeled collapse time is the first straight-line event after pulling
\end{definition}
\begin{proof}
  This is the declared model definition of Is First Straight Line Collapse.
\end{proof}
\begin{definition}[Answer Choice]
  \label{def:physics:phyx-mini-0884:phyxminiproblems-problemphyxmini0884-answerchoice}
  \lean{PhyXMiniProblems.ProblemPhyXMini0884.AnswerChoice}
  Labels of the four answer choices supplied with the problem
\end{definition}
\begin{proof}
  This is the declared model definition of Answer Choice.
\end{proof}
\begin{definition}[displayed Duration In Seconds]
  \label{def:physics:phyx-mini-0884:phyxminiproblems-problemphyxmini0884-answerchoice-displayeddurationinseconds}
  \lean{PhyXMiniProblems.ProblemPhyXMini0884.AnswerChoice.displayedDurationInSeconds}
  \uses{def:physics:phyx-mini-0884:phyxminiproblems-problemphyxmini0884-answerchoice}
  Duration in seconds printed beside each displayed answer choice
\end{definition}
\begin{proof}
  This is the declared model definition of displayed Duration In Seconds.
\end{proof}
\begin{definition}[recorded Dataset Answer]
  \label{def:physics:phyx-mini-0884:phyxminiproblems-problemphyxmini0884-recordeddatasetanswer}
  \lean{PhyXMiniProblems.ProblemPhyXMini0884.recordedDatasetAnswer}
  \uses{def:physics:phyx-mini-0884:phyxminiproblems-problemphyxmini0884-answerchoice}
  The answer label recorded in the supplied dataset metadata
\end{definition}
\begin{proof}
  This is the declared model definition of recorded Dataset Answer.
\end{proof}
\begin{definition}[Rounds To Nearest Tenth Second]
  \label{def:physics:phyx-mini-0884:phyxminiproblems-problemphyxmini0884-roundstonearesttenthsecond}
  \lean{PhyXMiniProblems.ProblemPhyXMini0884.RoundsToNearestTenthSecond}
  Rounding to the nearest tenth of a second, matching the answer display
\end{definition}
\begin{proof}
  This is the declared model definition of Rounds To Nearest Tenth Second.
\end{proof}
\begin{definition}[Answer Matches Collapse Time]
  \label{def:physics:phyx-mini-0884:phyxminiproblems-problemphyxmini0884-answermatchescollapsetime}
  \lean{PhyXMiniProblems.ProblemPhyXMini0884.AnswerMatchesCollapseTime}
  \uses{def:physics:phyx-mini-0884:phyxminiproblems-problemphyxmini0884-collapsingsquareloopsetup, def:physics:phyx-mini-0884:phyxminiproblems-problemphyxmini0884-collapseelapsedtimeinseconds, def:physics:phyx-mini-0884:phyxminiproblems-problemphyxmini0884-answerchoice, def:physics:phyx-mini-0884:phyxminiproblems-problemphyxmini0884-roundstonearesttenthsecond}
  A displayed choice agrees with the modeled physical collapse duration
\end{definition}
\begin{proof}
  This is the declared model definition of Answer Matches Collapse Time.
\end{proof}
\begin{lemma}[collapse elapsed time formula]
  \label{lem:physics:phyx-mini-0884:phyxminiproblems-problemphyxmini0884-collapse-elapsed-time-formula}
  \lean{PhyXMiniProblems.ProblemPhyXMini0884.collapse_elapsed_time_formula}
  \uses{def:physics:phyx-mini-0884:phyxminiproblems-problemphyxmini0884-speedinmeterspersecond, def:physics:phyx-mini-0884:phyxminiproblems-problemphyxmini0884-collapsingsquareloopsetup, def:physics:phyx-mini-0884:phyxminiproblems-problemphyxmini0884-sidelengthinmeters, def:physics:phyx-mini-0884:phyxminiproblems-problemphyxmini0884-collapseelapsedtimeinseconds, def:physics:phyx-mini-0884:phyxminiproblems-problemphyxmini0884-hasphysicalcollapsingloopparameters, def:physics:phyx-mini-0884:phyxminiproblems-problemphyxmini0884-satisfiesinextensiblesquaregeometry, def:physics:phyx-mini-0884:phyxminiproblems-problemphyxmini0884-satisfiesoppositecornerseparationkinematics, def:physics:phyx-mini-0884:phyxminiproblems-problemphyxmini0884-isfirststraightlinecollapse}
  The initial diagonal is √2 side lengths, the collapsed diagonal is two side lengths, and the diagonal opens at the sum of the two corner speeds. Hence the elapsed collapse time is the remaining diagonal distance divided by that relative speed. This is a conclusion, not a premise of the main theorem
\end{lemma}
\begin{proof}
  The conclusion follows from the governing relations and figure readouts named in the statement of collapse elapsed time formula.
\end{proof}
% --- Archon declaration topology end ---
% --- Archon physics formalization source end ---
